\section{敏感性分析}

\subsection{问题一覆盖宽度对坡度与开角的敏感性}

为考察几何模型对关键输入参数的敏感程度，在 $x=0$、$D_0=70\text{ m}$ 处分别单独扰动坡度 $\alpha$ 与换能器开角 $\theta$，覆盖宽度变化如表 5 所示。

\begin{table}[H]
  \centering
  \caption{问题一覆盖宽度对坡度与开角的单因素敏感性}
  \label{tab:sens}
  \begin{tabular}{lcccccc}
    \toprule
    $\alpha$ / $^\circ$（$\theta=120^\circ$） & 0.0 & 0.5 & 1.0 & 1.5 & 2.0 & 3.0 \\
    $W$ / m & 242.49 & 242.54 & 242.71 & 242.99 & 243.38 & 244.50 \\
    \midrule
    $\theta$ / $^\circ$（$\alpha=1.5^\circ$） & 100 & 110 & 120 & 130 & 140 & -- \\
    $W$ / m & 167.01 & 200.22 & 242.99 & 301.18 & 386.65 & -- \\
    \bottomrule
  \end{tabular}
\end{table}

由表 5 可知，覆盖宽度对坡度 $\alpha$ 的敏感程度较低，在 $0.0^\circ\sim3.0^\circ$ 范围内 $W$ 仅变化约 $2.01\text{ m}$；而覆盖宽度对开角 $\theta$ 明显更敏感，$\theta$ 从 $100^\circ$ 增至 $140^\circ$ 时 $W$ 由 $167.01\text{ m}$ 增至 $386.65\text{ m}$。这说明在本题给定的小坡度范围内，可把坡度近似作为弱影响因素，而换能器开角是决定覆盖宽度的主要参数。

\subsection{问题三最小测线数的约束敏感性}

问题三的递推求解取重叠率下限 $10\%$ 作为每步最大间距条件，方案恰好落在约束边界上。若在相同布线方向下减少一条测线，最东覆盖边界由 $3719.41\text{ m}$ 降至 $3678.24\text{ m}$，出现 $25.76\text{ m}$ 缺口。该结果表明 $34$ 条测线对该 $10\%\sim20\%$ 重叠率约束是紧的，方案对重叠率下限的变化较为敏感。

\subsection{问题四分带粒度的敏感性}

问题四的分带高度从 $5.00$ 海里细化到 $0.25$ 海里时，测线总长度由 $555600.00\text{ m}$ 降至 $415311.00\text{ m}$，超 $20\%$ 重叠长度由 $276877.69\text{ m}$ 降至 $2742.38\text{ m}$，漏测率始终为 $0.00\%$。可见分带粒度对冗余控制影响显著；推荐 $0.50$ 海里方案是在三项指标与线段数、转场成本之间的折中选择。
